%% Appendix A -- instrument and retrieval configuration (Phase 3, 2026-09-28).
%% Sources: operational fit-setting file
%%   C:\GHL\2026 yeosu\Output\fit setting\FitSet_ANs[430-462nm_P4]_PNs[444-471nm_P3]_cold[438-476nm_P4]_Std.json
%% (windows, orders, references, shift/squeeze settings, step limit, d, R_L, temperature coefficients),
%% the wavelength-calibration files it names (dispersion, FWHM of the heated channels), and the existing text
%% (cold FWHM, Sect. 4.1; R(t) knots, Sect. 6.4; calibration schedule, formerly Table 18).
\section{Instrument and retrieval configuration}
\label{app:A}

Augur reads the layout of the raw files -- which columns hold which channel's spectrum, which state flag marks an
ambient, zero-air or helium scan, and where the housekeeping values sit -- from a JSON instrument and campaign
profile, and the fit configuration from a separate fit-setting file. The retrieval logic keys on the structure these
files declare, not on channel names, so the same code is configured for another instrument or deployment by writing
new files rather than by changing it. This appendix collects the configuration of the field case of
Sect.~\ref{sec:7} in one place. Table~\ref{tab:config} is transcribed from the operational fit-setting file of the
campaign and the wavelength-calibration files it names; Table~\ref{tab:1} is its short form.

\begin{table*}[t]
\caption{Operational retrieval configuration of the three channels. ``Linked'' parameters take the value of the NO$_2$
parameter. The squeeze is stored as the departure from unity of $s_i$ in Eq.~(\ref{eq:3}), so 0 means no stretch.
Dispersion is the median over the fitting window of the wavelength-calibration file; the line-shape width of the
heated channels is the value recorded in that file (mean of four Hg lines), that of the ambient-temperature channel is
the pair of measurements discussed in Sect.~\ref{sec:4.1}.}
\label{tab:config}
\begin{tabular}{>{\raggedright\arraybackslash}p{0.27\textwidth}>{\raggedright\arraybackslash}p{0.2\textwidth}>{\raggedright\arraybackslash}p{0.2\textwidth}>{\raggedright\arraybackslash}p{0.22\textwidth}}
\tophline
 & 300\,$^{\circ}$C & 180\,$^{\circ}$C & ambient temperature \\
\middlehline
label in Sects.~\ref{sec:3}--\ref{sec:6} & hot ANs & hot PNs & cold \\
raw spectrum block & 2053 & 4101 & (cold file) \\
fitting window (nm) & 429.5--461.9 & 444.1--470.6 & 438.4--475.8 \\
baseline & Chebyshev, order 4 & Chebyshev, order 3 & Chebyshev, order 4 \\
absorbers fitted & NO$_2$, CHOCHO, H$_2$O & NO$_2$, CHOCHO, H$_2$O & NO$_2$, CHOCHO, H$_2$O \\
NO$_2$ shift (px) & bounded, $-10$ to $+0.5$ & bounded, $-5$ to $+5$ & fixed, $-0.5$ \\
NO$_2$ squeeze, $s-1$ & bounded, $\pm 0.005$ & bounded, $\pm 0.005$ & fixed, 0 \\
CHOCHO, H$_2$O shift and squeeze & linked & linked & linked \\
cross-section temperature coefficient $\kappa_i$ & 0 (not applied) & 0 (not applied) & 0 (not applied) \\
negative amounts & allowed & allowed & allowed \\
optimiser step limit & on (0.5) & on (0.5) & on (0.5) \\
wavelength calibration & Hg, 19 June 2026, quadratic & Hg, 19 June 2026, quadratic & Hg (four lines), 23 May 2026, quadratic \\
dispersion (nm\,px$^{-1}$) & 0.0484 & 0.0481 & 0.0482 \\
line-shape FWHM (nm) & 0.681 & 0.683 & 0.765--0.785 \\
\bottomhline
\end{tabular}
\end{table*}

\emph{References.} The cross sections are NO$_2$ \citep{Vandaele_2002}, CHOCHO \citep{Volkamer_2005} and H$_2$O from
HITRAN line data \citep{Gordon_2022}, each convolved with the line shape of its channel before the fit; O$_4$ is not
fitted in the operational configuration. No robust weighting, regularisation or non-negativity constraint is applied
(Appendix~\ref{app:B6}).

\emph{Wavelength calibration.} Each channel has its own Hg-lamp calibration with a quadratic pixel-to-wavelength
polynomial. Two calibrations of the 300\,$^{\circ}$C spectral region exist, one per region of interest on the
detector; the operational fit setting uses the second, the two differ by about 0.6\,px (0.034\,nm), and which of the
two belongs to block 2053 is not settled. The difference is within the shift range of that channel and is absorbed
by the fitted shift.

\emph{Temperature and pressure.} Sample temperature and pressure are recorded with every spectrum and enter the
Rayleigh correction of Eq.~(\ref{eq:1}) and the conversion from number density to mixing ratio; their interpolation
contributes below \fieldnum{0.1}\,\% (Sect.~\ref{sec:4.1}). The temperature scaling $\tau_i(T)$ of Eq.~(\ref{eq:2}) is not
used ($\kappa_i = 0$).

\emph{Cavity.} All three cavities have $d = 51.8$\,cm and $R_L = 1.0$. $R_L$ is assumed, not measured; since
$\alpha \propto R_L$ it enters as a multiplicative factor (Sect.~\ref{sec:3.6}).

\emph{Calibration references.} Zero-air and helium blocks are identical, 34 scans over 32\,s; zero air is measured
every hour and helium every three hours (Table~\ref{tab:cal_schedule}). The zero-air reference $I_0$ is interpolated
to the time of each spectrum from the neighbouring zero-air blocks, and the etalon frequency is held fixed
(Sect.~\ref{sec:4.1}). Each reflectivity knot is formed from one zero-air block and the most recent helium block that
passed the quality gates, so the knots of the heated channels are hourly -- \fieldnum{1261} over the deployment at a median
spacing of \fieldnum{1.00}\,h -- while the helium half of a knot is up to three hours old; $R(\lambda,t)$ is the monotone
piecewise-cubic interpolation of Sect.~\ref{sec:2.4} through these knots, after the knot gate of Sect.~\ref{sec:6.3}.

\begin{table}[t]
\caption{Calibration schedule of the heated channels, with block counts for the budget window (18--24 May): number of blocks, scans and duration per block, cadence and duty cycle, for zero air and helium.}
\label{tab:cal_schedule}
\begin{tabular}{lllll}
\tophline
 & blocks & per block & cadence & duty \\
\middlehline
zero air & \fieldnum{157} & 34 scans / 32\,s & 1.00\,h & 0.90\,\% \\
helium & \fieldnum{54} & 34 scans / 32\,s & 3.00\,h & 0.31\,\% \\
 & & & & \textbf{1.20\,\%} \\
\bottomhline
\end{tabular}
\end{table}
